\documentclass[11pt,a4paper]{article}

% Guia de practica 04: fundamentos de algoritmos (Bloque B)
\usepackage[utf8]{inputenc}
\usepackage[T1]{fontenc}
\usepackage[spanish,es-noshorthands,es-nodecimaldot,es-tabla]{babel}
\usepackage{lmodern}
\usepackage[margin=2.3cm]{geometry}
\usepackage{amsmath,amssymb}
\usepackage{xcolor}
\usepackage{enumitem}
\usepackage[most]{tcolorbox}
\usepackage{fancyhdr}
\usepackage{titlesec}
\usepackage{array,tabularx,booktabs}
\usepackage{listings}
\usepackage{xurl}
\usepackage{hyperref}
\usepackage{parskip}

\definecolor{morado}{RGB}{74,20,99}
\definecolor{moradoMedio}{RGB}{123,63,157}
\definecolor{moradoClaro}{RGB}{224,208,235}
\definecolor{tealProfundo}{RGB}{0,90,99}
\definecolor{tealClaro}{RGB}{204,232,232}
\definecolor{acento}{RGB}{198,110,38}
\definecolor{acentoClaro}{RGB}{247,226,201}
\definecolor{grafito}{RGB}{51,51,51}
\definecolor{grisClaro}{RGB}{244,246,245}
\definecolor{codigoFondo}{RGB}{248,247,251}

\hypersetup{colorlinks=true,linkcolor=morado,urlcolor=moradoMedio,
  pdftitle={Guia de practica 04: Fundamentos de algoritmos},
  pdfauthor={Programacion y Base de Datos - UNL},bookmarksopen=true}
\titleformat{\section}{\Large\bfseries\color{morado}}{\thesection.}{0.5em}{}
\titleformat{\subsection}{\large\bfseries\color{morado}}{\thesubsection}{0.5em}{}
\titlespacing*{\section}{0pt}{1.3em}{0.6em}
\pagestyle{fancy}
\fancyhf{}
\setlength{\headheight}{14pt}
\renewcommand{\headrulewidth}{0.6pt}
\fancyhead[L]{\small\color{grafito}Programacion y Base de Datos -- Unidad 1}
\fancyhead[R]{\small\color{grafito}Practica 04: Algoritmos}
\fancyfoot[C]{\small\thepage}
\setlist[itemize]{itemsep=2pt,topsep=3pt}
\setlist[enumerate]{itemsep=4pt,topsep=3pt}

\lstdefinestyle{pseudo}{
  basicstyle=\ttfamily\small\color{grafito},
  backgroundcolor=\color{codigoFondo},
  frame=leftline,framerule=3pt,rulecolor=\color{moradoMedio},
  xleftmargin=16pt,framexleftmargin=5pt,
  breaklines=true,columns=fullflexible,keepspaces=true,
  showstringspaces=false,tabsize=4,aboveskip=7pt,belowskip=7pt
}
\lstnewenvironment{pseudocodigo}{\lstset{style=pseudo}}{}
\newtcolorbox{objetivo}{colback=tealClaro!40,colframe=tealProfundo,
  arc=2pt,boxrule=0.8pt,left=8pt,right=8pt,top=6pt,bottom=6pt,breakable}
\newtcolorbox{atencion}{colback=acentoClaro!65,colframe=acento,
  arc=2pt,boxrule=0.8pt,left=8pt,right=8pt,top=6pt,bottom=6pt,breakable}
\newcommand{\code}[1]{\texttt{#1}}
\newcommand{\nbfile}{\code{03\_fundamentos\_algoritmos.ipynb}}

\begin{document}

\begin{titlepage}
\centering
\vspace*{0.8cm}
{\large Universidad Nacional de Loja\\[3pt]
Facultad Agropecuaria y de Recursos Naturales Renovables\\[3pt]
Carrera de Ingenieria Ambiental}\par
\vspace{1.5cm}
{\Large\color{moradoMedio}\bfseries GUIA DE PRACTICA N.\textsuperscript{o} 04}\par
\vspace{10pt}
{\Huge\bfseries\color{morado}Fundamentos de algoritmos}\par
\vspace{8pt}
{\Large\color{morado}Problema, algoritmo y programa -- Bloque B}\par
\vspace{1cm}
\begin{objetivo}
\textbf{Proposito.} Analizar problemas ambientales sencillos y representar sus soluciones mediante entrada, proceso y salida, pseudocodigo, diagramas de flujo y pruebas de escritorio, antes de escribir codigo.
\end{objetivo}
\vspace{1cm}
\begin{tabular}{rl}
\textbf{Asignatura:} & Programacion y Base de Datos (2.\textsuperscript{o} ciclo) \\
\textbf{Docente:} & Guillermo Chuncho Morocho \\
\textbf{Ciclo academico:} & Septiembre 2026 -- Febrero 2027 \\
\textbf{Cuaderno de trabajo:} & \nbfile{} \\
\textbf{Duracion estimada:} & 4 horas (2 en clase + 2 de trabajo autonomo) \\
\textbf{Modalidad:} & Practica guiada y trabajo individual \\
\end{tabular}
\vfill
{\small Loja, Ecuador}
\end{titlepage}

\tableofcontents
\newpage

\section{Datos informativos}
\begin{center}
\renewcommand{\arraystretch}{1.3}
\begin{tabularx}{\textwidth}{>{\bfseries\color{morado}}p{4.1cm}X}
\toprule
Unidad & Unidad 1 -- Introduccion a la programacion \\
Bloque & B: Fundamentos de algoritmos \\
Tema & Problema, algoritmo, programa, descomposicion y representacion \\
Herramienta & Cuaderno Jupyter/Google Colab: \nbfile{} \\
Contexto & Datos de lluvia, cobertura forestal, agua y biodiversidad \\
Conocimientos previos & Uso basico de celdas Markdown y lectura de tablas \\
\bottomrule
\end{tabularx}
\end{center}

\section{Objetivos}
\subsection{Objetivo general}
Diseñar y comprobar algoritmos para problemas ambientales sencillos, representandolos con el esquema entrada--proceso--salida, pseudocodigo y diagramas de flujo.

\subsection{Objetivos especificos}
\begin{enumerate}
  \item Diferenciar problema, algoritmo y programa con ejemplos de monitoreo ambiental.
  \item Aplicar descomposicion, abstraccion y reconocimiento de patrones para organizar un problema.
  \item Identificar entradas, procesos y salidas pertinentes para una solucion.
  \item Escribir pseudocodigo claro y elaborar un flujograma con simbolos adecuados.
  \item Realizar una prueba de escritorio y verificar resultados con casos limite.
\end{enumerate}

\section{Materiales y organizacion}
\begin{itemize}
  \item Computador o cuaderno de apuntes; acceso al cuaderno \nbfile{}.
  \item Lapiz, regla o herramienta de diagramacion para el flujograma.
  \item Datos de ejemplo incluidos en el notebook; no se requieren paquetes ni programacion en Python para esta practica.
\end{itemize}

\begin{atencion}
Esta guia trabaja el diseno del algoritmo antes de traducirlo a Python. No se evaluara que el estudiante escriba codigo: se evaluara que la solucion sea finita, ordenada, comprensible y comprobable.
\end{atencion}

\section{Fundamento teorico}
\subsection{Tres conceptos relacionados}
\begin{itemize}
  \item \textbf{Problema:} pregunta o situacion que se desea resolver.
  \item \textbf{Algoritmo:} secuencia finita y ordenada de pasos para obtener una solucion.
  \item \textbf{Programa:} algoritmo expresado en un lenguaje que la computadora puede ejecutar.
\end{itemize}
Un algoritmo puede revisarse en lenguaje cotidiano, pseudocodigo o flujograma sin depender todavia de Python.

\subsection{Pensamiento computacional}
\textbf{Descomponer} es dividir la tarea en partes manejables; \textbf{abstraer} es conservar lo relevante y dejar de lado detalles que no afectan la solucion; \textbf{reconocer patrones} es advertir pasos que se repiten. Estas herramientas ayudan a proponer un algoritmo antes de elegir instrucciones de un lenguaje.

\subsection{Entrada, proceso y salida}
\begin{center}
\renewcommand{\arraystretch}{1.35}
\begin{tabularx}{0.95\textwidth}{>{\bfseries\color{morado}}p{2.7cm}X}
\toprule
Entrada & Datos indispensables para resolver la pregunta. \\
Proceso & Operaciones, comparaciones y decisiones aplicadas a los datos. \\
Salida & Resultado interpretable que responde la pregunta, con contexto y unidades cuando corresponda. \\
\bottomrule
\end{tabularx}
\end{center}

\subsection{Pseudocodigo y flujograma}
El pseudocodigo describe instrucciones con lenguaje claro y estructura ordenada, sin exigir una sintaxis particular. En un flujograma, el ovalo representa inicio/fin, el paralelogramo entrada/salida, el rectangulo proceso, el rombo decision y las flechas indican el orden.

\subsection{Prueba de escritorio}
Consiste en simular a mano el algoritmo con datos pequenos, registrando como cambia cada variable o contador. Se recomienda probar un caso habitual y casos limite: una lista vacia, valores iguales al umbral o datos que no cumplen la condicion.

\section{Desarrollo de la practica}
Para cada actividad, primero lea la seccion correspondiente del notebook, redacte sus respuestas en una celda Markdown o en esta guia y compare su razonamiento con la prueba indicada. No copie una solucion sin poder explicar cada paso.

\subsection{Actividad 1. Del contexto a la pregunta}
\textbf{Consigna.} Un equipo dispone de registros diarios de precipitacion de una estacion y necesita informar cuantos dias del mes superaron los 10 mm.
\begin{enumerate}
  \item Redacte el problema como una pregunta concreta.
  \item Explique la diferencia entre el problema, el algoritmo que lo resuelve y el programa que podria ejecutarlo.
  \item Descomponga el trabajo en al menos tres subtareas ordenadas.
  \item Anote un dato necesario y un detalle irrelevante para responder la pregunta.
\end{enumerate}

\subsection{Actividad 2. Esquema entrada--proceso--salida}
Complete la tabla para el conteo de dias lluviosos. Sea preciso con el umbral: solo cuenta un dia si la lluvia es estrictamente mayor que 10 mm.
\begin{center}
\renewcommand{\arraystretch}{1.7}
\begin{tabularx}{\textwidth}{>{\bfseries\color{morado}}p{3cm}X}
\toprule
Entrada & \\
Proceso & \\
Salida & \\
\bottomrule
\end{tabularx}
\end{center}

\subsection{Actividad 3. Escribir pseudocodigo}
Escriba una secuencia finita de pasos que resuelva el conteo. Su propuesta debe leer la lista y el umbral, iniciar el contador, revisar cada observacion, actualizarlo solo cuando la condicion se cumpla y mostrar el resultado.

\textbf{Espacio de trabajo:}
\begin{pseudocodigo}
INICIO
    Escriba aqui su pseudocodigo
FIN
\end{pseudocodigo}

\textbf{Comprobacion:} identifique donde comienza el algoritmo, cual es su condicion, cuando cambia el contador y que evento detiene la revision de la lista.

\subsection{Actividad 4. Construir el flujograma}
Dibuje el algoritmo de la actividad 3. Incluya inicio y fin, lectura de datos, inicializacion del contador, decision sobre cada lluvia, actualizacion del contador, avance a la siguiente observacion y presentacion del resultado.
\begin{itemize}
  \item Use rombo para la comparacion \code{lluvia > 10} y rotule las ramas \textbf{Si} y \textbf{No}.
  \item Asegure que ambas ramas continuen correctamente y que el ciclo termine al revisar todos los datos.
  \item Revise que no existan pasos sin entrada o salida ni flechas ambiguas.
\end{itemize}

\subsection{Actividad 5. Prueba de escritorio}
Use las observaciones \code{[0, 12, 8, 25]} mm y umbral de 10 mm. Complete la tabla y determine el resultado final.
\begin{center}
\renewcommand{\arraystretch}{1.5}
\begin{tabularx}{0.85\textwidth}{c c c}
\toprule
Observacion (mm) & ¿Supera 10 mm? & Contador \\
\midrule
0 & & \\
12 & & \\
8 & & \\
25 & & \\
\bottomrule
\end{tabularx}
\end{center}
Responda: ¿que resultado espera? ¿Que ocurriria si una medicion fuera exactamente 10 mm? ¿Como cambiaria el resultado si el problema pidiera lluvia \emph{mayor o igual} a 10 mm?

\subsection{Actividad 6. Aplicacion: cobertura de bosque nativo}
Un equipo analiza cinco celdas con porcentajes de cobertura \code{[65, 40, 52, 50, 80]}. Debe informar cuantas tienen cobertura estrictamente superior al 50\,\% y que porcentaje representan del total.
\begin{enumerate}
  \item Defina el problema y las entradas.
  \item Describa el proceso y la salida usando E-P-S.
  \item Escriba pseudocodigo, incluyendo el conteo y el calculo del porcentaje.
  \item Haga una prueba de escritorio que muestre el contador despues de cada celda.
  \item Verifique manualmente el porcentaje final.
\end{enumerate}

\subsection{Actividad 7. Transferencia: calidad del agua}
Un laboratorio tiene mediciones de pH de varios puntos de una quebrada y necesita reportar su promedio. Sin programar, indique el problema, entradas, proceso y salida. Explique como modificaria el algoritmo si tambien se solicitara el numero de mediciones superiores a un valor de referencia entregado por el laboratorio.

\section{Buenas practicas de diseno}
\begin{itemize}
  \item Escriba pasos en orden y use verbos concretos: leer, comparar, acumular, calcular y mostrar.
  \item Defina los limites con precision: \code{>} no significa lo mismo que \code{>=}.
  \item Asegure que los nombres de las salidas expliquen que se conto o calculo.
  \item Compruebe que todo ciclo tiene una condicion de termino y que cada observacion se procesa una sola vez.
  \item Use datos pequenos para validar el razonamiento antes de traducirlo a un lenguaje.
\end{itemize}

\section{Autoevaluacion}
Responda con sus propias palabras:
\begin{enumerate}
  \item ¿En que se diferencian problema, algoritmo y programa?
  \item ¿Que datos son esenciales y cuales se pueden abstraer en un problema?
  \item ¿Que representa cada etapa de entrada, proceso y salida?
  \item ¿Que simbolos del flujograma representan una decision y un proceso?
  \item ¿Que permite detectar una prueba de escritorio antes de escribir codigo?
  \item ¿Por que es importante probar un valor igual al umbral?
\end{enumerate}

\section{Entregable y evaluacion}
\subsection{¿Que entregar?}
Entregue el notebook \code{P04\_Apellido\_Nombre.ipynb}, con las respuestas de las actividades 1 a 7 y la autoevaluacion. El cuaderno debe incluir el esquema E-P-S, pseudocodigo, flujograma (insertado como imagen o elaborado en una celda de texto) y las pruebas de escritorio. Puede exportarlo como \code{.ipynb} desde Colab/Jupyter o compartir el enlace segun las indicaciones del docente.

Antes de enviar, revise que los diagramas se lean con claridad y que el pseudocodigo coincida con la prueba de escritorio. No se requiere ejecutar codigo Python para esta guia.

\subsection{Rubrica}
\begin{center}
\renewcommand{\arraystretch}{1.3}
\begin{tabularx}{\textwidth}{>{\raggedright\arraybackslash}p{3.3cm}XXX}
\toprule
\textbf{Criterio} & \textbf{Logrado} & \textbf{En progreso} & \textbf{Por fortalecer} \\
\midrule
Analisis y E-P-S & Identifica problema y datos pertinentes; proceso y salida responden a la pregunta. & Presenta una omision menor. & Confunde etapas o selecciona datos sin relacion. \\
Pseudocodigo & Pasos claros, ordenados, finitos y coherentes con el criterio. & Solucion casi completa con alguna ambiguedad. & Pasos incompletos o no resuelven la pregunta. \\
Flujograma & Simbolos y rutas correctos; ciclo y fin son claros. & Se entiende con pequenos problemas de notacion. & Decisiones o flechas no representan el algoritmo. \\
Prueba y validacion & Tabla completa y resultado verificado manualmente. & Presenta un error menor de seguimiento. & No comprueba el contador o el resultado. \\
Comunicacion & Usa nombres, unidades y explicaciones precisas. & La explicacion es comprensible pero incompleta. & No puede explicar su representacion. \\
\bottomrule
\end{tabularx}
\end{center}

\section{Referencias}
\begin{itemize}
  \item Python Software Foundation. (2025). \emph{El tutorial de Python}. \url{https://docs.python.org/es/3/tutorial/}
  \item Downey, A. B. (2024). \emph{Think Python: How to think like a computer scientist} (3.\textsuperscript{a} ed.). \url{https://allendowney.github.io/ThinkPython/}
\end{itemize}

\end{document}
